%==============================================================================
\section{Background and Related Work}
\label{sec:background}
%==============================================================================

\textbf{RAG architectures.} RAG has developed from retrieval-augmented language models \cite{guu2020realm} and retrieve-then-generate pipelines \cite{izacard2021fid} into modular systems \cite{gao2024retrieval}. Modular RAG separates the retriever, generator, and evaluator into components that can be changed independently. Self-RAG \cite{asai2024selfrag} adds reflection tokens for retrieval and critique, but it is unclear whether self-correction can resist adversarial inputs designed to manipulate the critique process.

\textbf{Evaluation.} Retrieval quality is often measured with IR metrics such as top-$k$ retrieval accuracy \cite{karpukhin2020dense} and MRR. These metrics measure relevance, but not safety. A poisoned document that is highly relevant to a query can still score well. Generation quality has moved from n-gram overlap toward factuality-based measures that separate faithfulness from factual correctness \cite{maynez2020faithfulness}. TruthfulQA \cite{lin2022truthfulqa} focuses on imitative falsehoods, and the TRUE benchmark \cite{honovich2022true} supports factual-consistency evaluation with metrics such as NLI. Recent hallucination detectors mainly verify generated text rather than the integrity of the corpus. SelfCheckGPT \cite{manakul2023selfcheckgpt} measures self-consistency across sampled generations, while fine-grained factuality scores \cite{min2023factscore} check atomic facts against a trusted knowledge source. RAGAS \cite{es2024ragas} and ARES \cite{saadfalcon2024ares} use LLMs to judge faithfulness and relevance, but they do not evaluate whether the retrieved corpus itself is trustworthy.

\textbf{Attacks and defenses.} Indirect prompt injection places malicious instructions inside retrievable content \cite{greshake2023indirect}, and poisoning the web-scale datasets used for model pre-training can be done at low cost \cite{carlini2024poisoning}. PoisonedRAG \cite{zou2024poisoned} studies knowledge poisoning in the inference corpus, and adversarial passages can be optimized so that they enter the retrieved set for many queries \cite{zhong2023poisoning}. On the defense side, retrieval augmentation can reduce hallucination \cite{shuster2021retrieval}, and NLI-based verification can compare a retrieved document as the premise with the generated answer as the hypothesis to estimate factual consistency \cite{maynez2020faithfulness,honovich2022true}. However, many poisoning defenses are still \emph{offline}, such as filtering or cleaning the corpus before indexing. These methods cannot detect poisoned content that appears at inference time. Toxicity classifiers also miss misinformation that is written in neutral and professional language. Our work combines factual verification and poison detection into one interpretable online score, and shows where this type of surface-signal detection works and where it fails.

\textbf{Security in the software-development lifecycle.} Software-engineering research increasingly studies how generative AI is used in development \cite{tuomisto2025ethical} and how AI is adopted in quality assurance \cite{karhu2025barriers}, and how digital experimentation supports sustainability in software-intensive industry \cite{jantti2025digital}. LLM coding assistants are now widely used during development, but they can generate insecure code, and developers may over-trust their suggestions \cite{pearce2022,perry2023}. The OWASP Top 10 for LLM Applications lists data and model poisoning and misinformation as major risks for these systems \cite{owaspllm2025}. When a coding assistant uses RAG over an organization's secure-coding guidance, such as the OWASP Top 10 \cite{owasp2021} and CWE \cite{mitrecwe}, that guidance corpus becomes a poisoning target. A single corrupted rule can lead the assistant to give unsafe recommendations. This motivates evaluating the Evaluation Agent in an SDLC setting, \hyperref[sec:usecase]{Section~\ref*{sec:usecase}: Use Case: Secure-Coding RAG Assistant in the SDLC}.

